\begin{table}[!htbp]
  \centering
  \small
  \caption{Relationship between a representative concurrent-learning workflow
  and the validation-driven procedure used here.}
  \label{tab:si_concurrent_learning}
  \setlength{\tabcolsep}{4pt}
  \begin{tabular}{>{\raggedright\arraybackslash}p{2.0cm}>{\raggedright\arraybackslash}p{5.0cm}>{\raggedright\arraybackslash}p{7.1cm}}
    \toprule
    Aspect & DP-GEN concurrent learning & Present study \\
    \midrule
    Potential model & Deep Potential ensemble & Fine-tuned MACE-MPA-0 models \\
    Exploration & Iterative configuration-space sampling & Finite-temperature
    trajectories used to generate displaced test configurations \\
    Selection signal & Quantitative ensemble model deviation & High displacement
    and qualitative seed sensitivity; no calibrated deviation threshold \\
    First-principles step & Labels selected uncertain configurations & Tests MACE
    forces and identifies chemistry-specific failure modes \\
    Retraining loop & Selected labels are added and the cycle is repeated & The DFT
    snapshots are not added to a subsequent training cycle \\
    Primary purpose & Expand the range described by the potential & Determine
    which adsorption and force predictions are supported by DFT \\
    \bottomrule
  \end{tabular}
\end{table}
